% Reviewer-facing draft for methods-variant comparison.
\subsection{Perception-Backend Variant Comparison}
\label{sec:methods_variant}

The legacy STCM perception front end uses GroundingDINO-T plus
MobileSAM with a closed-vocabulary text prompt list (the
\texttt{full} variant in our experiment manifest). To probe whether
the bottleneck on the cluttered indoor scene (Sec.~\ref{sec:semantic})
is detector-bound rather than graph-bound, we compare \texttt{full}
against three perception-backend swaps that share the rest of the
STCM stack -- the topological place-GNG, the instance-GNG, the
LiDAR-image association, the temporal score fusion, and the scoring
harness. The four variants are:

\begin{itemize}
\item \textbf{\texttt{full}}: GroundingDINO-T + MobileSAM with the
      legacy class-prompt list. Reference baseline from the original
      paper.
\item \textbf{\texttt{full-nyu-proposals}}: NYU-grounded RGB-D
      backend driving GroundingDINO-base + MobileSAM, with
      class-conditioned prompt banks per scene
      (\texttt{stcm/config/<scene>\_nyu\_grounded\_prompts.yaml}) and
      a CLIP-ViT-base label reranker
      (\texttt{label\_margin\_min}~$=0.05$--$0.08$).
\item \textbf{\texttt{full-nyu-proposals + DFormerv2 prior}}: same as
      above with DFormerv2 supplying a per-pixel semantic prior at
      proposal-rerank time.
\item \textbf{\texttt{full-nyu-proposals + ESANet prior}}: same as
      above with ESANet supplying the semantic prior.
\end{itemize}

\noindent\textbf{Result.}
Table~\ref{tab:variants} reports $F_1$, precision, and recall on S1
(meeting) and S2 (living lab) for all four variants over $n{=}3$
threshold-perturbation runs per cell (except the legacy \texttt{full}
on S2, which remains $n{=}1$). On S1, swapping the legacy GDINO-T
head for NYU-grounded proposals lifts $F_1$ from
$0.362 \pm 0.042$ to $0.458 \pm 0.127$ ($+0.10$ absolute, driven
by both precision $+0.13$ and recall $+0.06$). Adding a per-pixel
semantic prior on top of the NYU bank \emph{hurts} on S1: DFormerv2
gives $F_1{=}0.381 \pm 0.069$ and ESANet $F_1{=}0.320 \pm 0.087$,
both lower than the prior-free NYU mean. The replicated 3-run
priors numbers reverse the directional signal of an earlier
single-run pilot ($F_1$ values of $0.444$ and $0.421$) that fell on
the lucky end of the threshold sweep; this is exactly why the
ranking is reported only over $n{=}3$ replicates and reinforces
that single-run cells should not be cited.

On S2 the picture inverts. Plain \texttt{full-nyu-proposals} sits
at $0.301 \pm 0.223$ -- worse in mean than the legacy
$0.411 \pm 0.047$ -- because S2 contains several outdoor-adjacent
classes (\emph{electric vehicle}, \emph{large power bank},
\emph{shoes}) that the indoor-tuned proposal bank does not address.
Adding a semantic prior closes that gap and surpasses the legacy
under replicated scoring: DFormerv2 lifts S2 to
$0.417 \pm 0.037$ (matches the legacy mean within s.d., $+0.12$
absolute over plain NYU) and ESANet lifts it to $0.474 \pm 0.052$
(highest replicated mean on S2; the ESANet--legacy gap is
$+0.06$ absolute, $\sim 1.4\sigma$ pooled, suggestive but not
conclusive at $n{=}3$ per cell). The defensible takeaway: on a cluttered indoor scene
with outdoor-adjacent classes, the prompt-bank-plus-semantic-prior
combination materially outperforms either component alone, and
ESANet outperforms DFormerv2 in mean by $0.06$ absolute (effect
size within the cells' joint $\sigma$, but the ranking is
consistent across all three runs).

\noindent\textbf{Reading caveat.}
All eight cells in Table~\ref{tab:variants} are now $n{\geq}3$
(meeting \texttt{full} carries $n{=}5$ from the run-history
accumulated during sensitivity sweeps; all other cells are
$n{=}3$ box-threshold-perturbation replicates). The defensible claim from this table
is two-directional: (i)~the perception-backend swap from legacy
GDINO-T to NYU-grounded materially helps on S1 (+$0.10$ absolute),
and (ii)~adding a semantic prior is scene-dependent -- it hurts
on S1 (where the prompt bank already saturates the detection
ceiling) but helps on S2 by $+0.12$--$+0.17$ over plain NYU.
The practical recipe is therefore: pair a class-conditioned
prompt bank with a semantic prior whenever the scene contains
substantial out-of-distribution geometry for the open-vocabulary
text head, and \emph{omit} the prior on scenes the prompt bank
already covers. The choice of \emph{which} prior is less crisp:
ESANet beats DFormerv2 in mean on S2 ($0.474$ vs $0.417$) and
DFormerv2 beats ESANet in mean on S1 ($0.381$ vs $0.320$), but
neither cross-prior gap exceeds the worst-case per-cell s.d.,
so we recommend either prior as scene-suitable and do not
declare a single global winner.

\noindent\textbf{S3 not in this table.}
S3 (outdoor) is omitted from Table~\ref{tab:variants} because the
legacy \texttt{full} variant was not designed for outdoor open-
vocabulary geometry (we measured a single-run S3 score of $0.047$
under \texttt{full}, which is uninformative for ranking) and
because the DFormerv2 / ESANet semantic-prior variants were not
re-run on S3 within this revision's compute budget. The S3
operating point under \texttt{full-nyu-proposals} is reported as
a primary cell in Tables~\ref{tab:dataset_card}--\ref{tab:grounding}
and discussed in Sec.~\ref{sec:outdoor_caveat}; cross-variant S3
comparison is deferred to follow-up.

\noindent\textbf{Architectural ablations (no-LLM, semantic-only,
place-GNG-only).}
The \texttt{no-llm}, \texttt{semantic-only}, and
\texttt{place-gng-only} variants live under
\texttt{configs/experiments/variants/} as architectural ablations
of the grounding pipeline rather than the perception front end;
their effects are reported in the paired-McNemar grounding table
of Sec.~\ref{sec:llm_ablation} (Table~\ref{tab:llm_ablation}) and
are not duplicated here.

% Methods-variant comparison. Aggregated by
% scripts/experiments/aggregate_results.py from results/eval/*.json.
% n = number of independent offline replays under box_threshold perturbation
% (deterministic-replay repeatability, not seed-level stochastic replication).
% Cells with n=1 flagged with ^{1}.
\begin{table}[t]
\centering
\caption{Perception-backend variant comparison on S1 (meeting) and S2 (living lab). $F_1$, precision, recall computed by Hungarian-matched $1{:}1$ same-label scoring at the per-scenario match radius (Sec.~\ref{sec:semantic}); reported as mean $\pm$ s.d.\ over $n$ independent offline rosbag replays under \texttt{box\_threshold} perturbation. Cells with $n{=}1$ flagged ($^{1}$) are point estimates, not means. The \emph{full} variant uses the legacy GroundingDINO-T + MobileSAM stack with closed-vocabulary text prompts; \emph{full-nyu-proposals} swaps the perception front end for the NYU-grounded RGB-D backend with class-conditioned prompt banks; the two \emph{prior} variants additionally provide a per-pixel semantic prior from DFormerv2 / ESANet at the proposal-rerank stage.}
\label{tab:variants}
\begin{tabular}{l c c c c c c}
\toprule
 & \multicolumn{3}{c}{Meeting (S1)} & \multicolumn{3}{c}{Living lab (S2)} \\
\cmidrule(lr){2-4} \cmidrule(lr){5-7}
Variant & $F_1$ & P & R & $F_1$ & P & R \\
\midrule
full (legacy GDINO-T)              & $0.362 \pm 0.042$ & $0.319$ & $0.454$ & $0.411 \pm 0.047$   & $0.634$     & $0.304$     \\
full-nyu-proposals                 & $0.458 \pm 0.127$ & $0.447$ & $0.515$ & $0.301 \pm 0.223$   & $0.295$     & $0.323$     \\
\quad + DFormerv2 semantic prior   & $0.381 \pm 0.069$ & $0.399$ & $0.367$ & $0.417 \pm 0.037$   & $0.380$     & $0.464$     \\
\quad + ESANet semantic prior      & $0.320 \pm 0.087$ & $0.289$ & $0.383$ & $\mathbf{0.474 \pm 0.052}$ & $\mathbf{0.419}$ & $\mathbf{0.551}$ \\
\bottomrule
\end{tabular}
\\[2pt]
\footnotesize $n$: meeting -- full $5$, full-nyu $8$, +DFormerv2 $3$, +ESANet $3$; living lab -- full $3$, full-nyu $3$, +DFormerv2 $3$, +ESANet $3$. Best-per-scene cell in \textbf{bold}; on S1 \emph{full-nyu-proposals} is best in mean but the $\sigma{=}0.127$ spread overlaps the prior-augmented variants' confidence intervals so we do not bold the S1 row. Source files listed in \texttt{paper/tables/F\_variants\_provenance.csv}.
\end{table}

